\newpage
	
	\pagestyle{plain}
	\setcounter{page}{1}
	
	\chapter{Удиртгал}
	\section{Үндэслэл, ач холбогдол}
	
	Өнөө үед дижитал технологи залуу үеийн өдөр тутмын амьдралын салшгүй хэсэг болж, мэдлэг мэдээллээ олж авах үндсэн сувгуудын нэг болсон. Ухаалаг гар утас, таблет болон онлайнаар дамжих контентууд уламжлалт сурах бичиг, хэвлэмэл номоос хамаагүй түгээмэл хэрэглэгддэг болжээ. Ялангуяа өсвөр үе болон оюутан залуучуудын хувьд мэдээлэл авах хэлбэр цахим хэрэглээ рүү шилжсэн ч, тэдэнд зориулсан чанартай, системтэй, интерактив боловсролын аппликэйшн Монгол хэл дээр хангалттай түгээмэл бус байна.

Монголын түүх нь дэлхийн түүхэнд онцгой байр суурь эзэлдэг хэр нь залуучуудад ойлгомжтой, сонирхол татахуйц, дижитал хэлбэрээр хүрэх боломж хомс хэвээр байна. Одоогоор Монгол хэл дээрх боловсролын апп-ууд цөөн, контентын чанар болон интерактив байдал сул, хэрэглэгчдийг удаан хугацаанд татах боломж хязгаарлагдмал байдаг. Үүнээс шалтгаалж залуус түүхийн мэдлэгээ гүнзгийрүүлэхээс илүү богино хэмжээний, түр зуурын анхаарал татсан контентууд руу ханддаг болсон.

Тиймээс энэхүү төсөл нь зөвхөн нэгэн апп бүтээх зорилготой бус, харин Монголын түүхийг орчин үеийн технологийн шийдлээр баяжуулж, залуу үедээ хүртээмжтэй, сонирхолтой хэлбэрээр хүргэхэд чиглэнэ. Интерактив визуал контент, тоглоомын элемэнт, түвшин ахих систем зэрэг нь хэрэглэгчдийн сурах сонирхлыг нэмэгдүүлж, түүхийн талаарх мэдлэг, үндэсний ухамсрыг бодитоор дэмжих боломжтой.

Ингэснээр уг систем нь боловсролын салбарт дижитал шинэчлэл хийх, Монголын түүхийг илүү өргөн хүрээнд түгээх, залуу үеийн түүхийн мэдлэгийг нэмэгдүүлэх зэрэг олон талын ач холбогдолтой гэж үзэж байна.
	
	\section{Зорилго, зорилт}
	
	\subsection*{Зорилго}
	
	13-р зууны Монголын түүхийг таниулах интерактив апп бүтээх замаар залуу үеийнхэнд түүхийн мэдлэгийг түгээх, сонирхлыг нэмэгдүүлэх. Энэ нь түүхийн боловсролыг дижитал хэлбэрээр хүртээмжтэй болгох, Монгол хэл дээрх боловсролын нөөцийн дутагдлыг нөхөхөд чиглэнэ.
	
	\newpage
	\subsection*{Зорилт}
	
	\begin{enumerate}
		\item Түүхийн эх сурвалж судлах, мэдээлэл цуглуулах.
		\item Апп-ын шаардлага, дизайныг тодорхойлох (SRS, UI/UX, ERD).
		\item Flutter ашиглан кросс-платформ апп бүтээх.
		\item SQLite-д өгөгдөл хадгалах, Firebase-ээр онлайн sync хийх.
		\item Туршилт явуулж, хэрэглэгчийн санал хүсэлтээр сайжруулах.
		\item Үр дүнг дүгнэж, цаашдын өргөтгөлийг төлөвлөх.
	\end{enumerate}

	\section{Системийн функцууд}
	
	Апп нь дараах үндсэн функцуудыг агуулна:
	
	\begin{table}[h]
	\centering
	\begin{tabular}{|p{1.5cm}|p{5cm}|p{7cm}|}
	\hline
	\textbf{Код} & \textbf{Функц} & \textbf{Тайлбар} \\
	\hline
	FR-01 & Түүхэн хүмүүсийн намтар & Чингис хаан, Өгэдэй зэрэг хүмүүсийн намтар, зураг харуулах \\
	\hline
	FR-02 & Байлдан дагуулалтын газрын зураг & Интерактив газрын зураг, зум хийх, тулаан дээр дарах \\
	\hline
	FR-03 & Цагийн хэлхээс & Үйл явдлуудыг цаг хугацаагаар харуулах timeline \\
	\hline
	FR-04 & Квиз & Мэдлэг шалгах асуулт хариулт, оноо цуглуулах \\
	\hline
	FR-05 & Мультимедиа & Зураг, видео, аудио үзэх \\
	\hline
	FR-06 & Гэр бүлийн мод & Хаадын гэр бүлийн мод харуулах \\
	\hline
	FR-07 & Соёлын мэдээлэл & Соёл, нийгмийн амьдралын мэдээлэл \\
	\hline
	\end{tabular}
	\caption{Системийн функцуудын жагсаалт}
	\label{tab:functions}
	\end{table}

	\section{Хэрэглэгчийн зорилтот бүлэг}
	
	\begin{itemize}
		\item \textbf{Үндсэн хэрэглэгч:} Залуучууд, сурагчид, оюутнууд (13+ насны)
		\item \textbf{Нэмэлт хэрэглэгч:} Багш нар, түүх судлаачид
	\end{itemize}
	
	Хэрэглэгчид нь Монгол хэл мэдэх, Android эсвэл iOS мобайл төхөөрөмжтэй байх шаардлагатай.
